\chapter{Qué resuelve Alexandria y con qué piezas}
\label{ch:modelo}

\section{El problema: una pila de especificaciones con huecos}

Un espacio de datos no es una especificación: es una \textbf{pila} de ellas. ODRL para las políticas, DCAT para el
catálogo, DSP para el protocolo, DCP para las credenciales. Cada una deja libertad a propósito. ODRL define qué es
una restricción, pero no qué operandos existen en tu sector. DSP define un \cod{ContractRequestMessage}, pero no
qué política es aceptable dentro.

Cada una de esas libertades es un \textbf{slot}. Un espacio de datos concreto lo cierra con un \textbf{binding}:
qué perfil manda en ese slot, en qué versión, con qué obligación y en qué punto del tráfico se comprueba, que es
el \textbf{gate}.

\begin{ideaClave}
La autoridad del espacio escribe esa pila \emph{una sola vez}, en YAML. Todo lo demás (ontología, contexto JSON-LD,
JSON Schema, \emph{shapes} SHACL, especificación legible) se \textbf{genera}. Nadie edita a mano dos artefactos
que deban decir lo mismo, porque acaban diciendo cosas distintas.
\end{ideaClave}

El problema no termina al generar. Los participantes tienen que enterarse de que hay una versión nueva, fiarse de
que viene de quien dice venir, aplicarla todos a la vez y no romper los contratos ya firmados cuando una regla
cambia. De eso se ocupan el hub y el cliente.

\section{Las tres piezas}

\begin{center}
\begin{tikzpicture}[
  node distance=9mm and 19mm, font=\small,
  box/.style={draw=alxNavy, fill=alxLightBlue, rounded corners=2pt, minimum height=11mm, minimum width=30mm, align=center},
  ext/.style={draw=alxBorder, fill=alxSlate, rounded corners=2pt, minimum height=11mm, minimum width=24mm, align=center},
  art/.style={draw=alxTeal, fill=alxLightTeal, rounded corners=2pt, minimum height=9mm, minimum width=30mm, align=center},
  arr/.style={-{Stealth[length=2.2mm]}, alxNavy, thick}]
  \node[ext] (yaml) {Bundle YAML\\[-1pt]\scriptsize perfil + espacio};
  \node[box, right=of yaml] (hub) {\texttt{alexandria-server}\\[-1pt]\scriptsize hub};
  \node[art, right=of hub] (pub) {Paquete firmado\\[-1pt]\scriptsize \texttt{/\allowbreak{}pub}, listado, registro};
  \node[box, below=14mm of pub] (cli) {\texttt{alexandria-client}\\[-1pt]\scriptsize borde L7};
  \node[ext, left=of cli] (conn) {Conector DSP\\[-1pt]\scriptsize sin modificar};
  \node[ext, right=of cli] (peer) {Contraparte};
  \node[box, below=14mm of yaml] (cc) {\texttt{alexandriac}\\[-1pt]\scriptsize compilador offline};
  \draw[arr] (yaml) -- node[above]{\scriptsize borrador} (hub);
  \draw[arr] (hub) -- node[above]{\scriptsize aprueba} (pub);
  \draw[arr] (pub) -- node[right]{\scriptsize sincroniza y verifica} (cli);
  \draw[arr] (peer) -- node[above]{\scriptsize mensaje} (cli);
  \draw[arr] (cli) -- node[above]{\scriptsize reenvía} (conn);
  \draw[arr, dashed] (yaml) -- node[left]{\scriptsize mismo código} (cc);
\end{tikzpicture}
\end{center}

\noindent\begin{tabularx}{\linewidth}{@{}L{3.5cm}L{2.7cm}Y@{}}
\toprule
\textbf{Binario} & \textbf{Quién lo usa} & \textbf{Qué hace} \\
\midrule
\cod{alexandriac} & Quien escribe el perfil & Valida y compila YAML en local, sin red y sin estado. Es el bucle rápido de edición. \\
\cod{alexandria-server} & La autoridad del espacio & Gobierna el ciclo borrador, candidato, versión y revocación. Firma lo que publica y lo sirve inmutable bajo \cod{/pub}. \\
\cod{alexandria-client} & Cada participante & Se coloca delante del conector. Sincroniza las versiones, verifica las firmas y decide sobre cada mensaje con estado local. \\
\bottomrule
\end{tabularx}

\medskip
El hub usa el mismo paquete Go que \cod{alexandriac} (\cod{internal/dsl}). Lo que compila en tu portátil compila
igual en el hub.

Hay además dos binarios de laboratorio que no son producto: \cod{fake-idp}, un emisor OAuth/OIDC que firma
tokens, y \cod{fake-connector}, un conector DSP de mentira que responde \cod{200} a todo y recuerda lo que recibe.

\section{El caso de estudio: Flexgrid}

Flexgrid es un espacio de datos sectorial ficticio. Los participantes y sus intereses son estos:

\noindent\begin{tabularx}{\linewidth}{@{}L{3.3cm}Y@{}}
\toprule
\textbf{Participante} & \textbf{Qué quiere} \\
\midrule
Distribuidora (DSO) & Publica la curva de carga cuartohoraria de sus puntos de suministro. Es la proveedora y quiere limitar para qué se usa. \\
Agregador & Entrena modelos de previsión de demanda sobre datos agregados por zona. \\
Comercializadora & Liquida desvíos y factura. Hasta ahora también perfilaba consumidores individuales. \\
Universidad & Quiere los datos para investigación, una finalidad que al principio no existe. \\
Autoridad del espacio & Decide el vocabulario común y lo publica. Lucía, de la distribuidora, actúa como \emph{steward}. \\
\bottomrule
\end{tabularx}

\medskip
Las reglas de negocio que la autoridad quiere poder expresar, y con qué pieza de ODRL se corresponden:

\noindent\begin{tabularx}{\linewidth}{@{}Y L{2.6cm} L{4.2cm}@{}}
\toprule
\textbf{Regla de negocio} & \textbf{Pieza ODRL} & \textbf{Término} \\
\midrule
«Puedes usarlos para prever la demanda.» & Acción & \cod{enr:forecastDemand} \\
«Puedes usarlos para liquidar desvíos.» & Acción & \cod{enr:settleImbalance} \\
«Solo para esta finalidad.» & Operando izquierdo & \cod{odrl:purpose}, adoptado y acotado \\
«Cada registro agrega al menos 50 puntos de suministro.» & Operando izquierdo & \cod{enr:minAggregation} \\
«No los guardes más de 90 días.» & Operando izquierdo & \cod{enr:retentionDays} \\
«Solo se tratan en España o en la UE.» & Operando izquierdo & \cod{enr:processingRegion} \\
«El permiso caduca el 1 de enero.» & Operando izquierdo & \cod{odrl:dateTime}, adoptado \\
\bottomrule
\end{tabularx}

\medskip
Sobre este espacio se cuenta además una historia de gobierno, que es la que da sentido al ciclo de vida:

\begin{enumerate}[leftmargin=*, itemsep=2pt]
  \item \textbf{1.0.0.} Primera versión. Incluye la acción \cod{enr:profileConsumer}, perfilar a un consumidor individual.
  \item \textbf{1.1.0.} La universidad entra: se añade la finalidad \cod{enr:research}. A la vez el regulador avisa de
        que el perfilado individual no va a ser admisible, y la acción se marca obsoleta. Cambio \emph{compatible}.
  \item \textbf{2.0.0.} El perfilado se retira. Cambio de \emph{ruptura}, con una ventana de gracia para migrar.
  \item \textbf{Revocación.} El regulador declara ilícito el perfilado. La autoridad revoca la 1.1.0, y los acuerdos
        firmados bajo ella dejan de valer para nuevas transferencias.
\end{enumerate}

\section{Preparación}

Desde la raíz del repositorio:

\begin{lstlisting}[style=sh]
make build
\end{lstlisting}
\begin{lstlisting}[style=out]
go build -o bin/alexandriac ./cmd/alexandriac
go build -o bin/alexandria-server ./cmd/alexandria-server
go build -o bin/alexandria-client ./cmd/alexandria-client
go build -o bin/fake-idp ./deploy/fake-idp
go build -o bin/fake-connector ./deploy/fake-connector
\end{lstlisting}

Los capítulos de la parte II solo necesitan \cod{bin/alexandriac}. Los de las partes III y IV usan los cinco.
